\documentclass[11pt]{article}
% preamble.tex — shared preamble for the Flyxion connective-essay series
% Compile target: XeLaTeX. \input this file, or copy its contents, at the
% top of each essay after \documentclass.

\usepackage{fontspec}
\setmainfont{TeX Gyre Termes} % Times-like serif; falls back gracefully if absent
\usepackage{amsmath,amssymb,amsthm}
\usepackage[margin=1.25in]{geometry}
\usepackage[hidelinks]{hyperref}
\usepackage{microtype}
\usepackage{enumitem}
\usepackage{footmisc}

\newtheorem{claim}{Claim}
\newtheorem{definition}{Definition}
\newtheorem{observation}{Observation}

% --- Corpus vocabulary macros -------------------------------------------
\newcommand{\admissible}{\textsc{admissible}}
\newcommand{\inadmissible}{\textsc{inadmissible}}
\newcommand{\repair}{\textit{repair}}
\newcommand{\continuation}{\textit{continuation}}
\newcommand{\rsvp}{\textsc{rsvp}}

% --- Title block helper ---------------------------------------------------
% Usage: \essaytitle{Title}{Subtitle or tagline}
\newcommand{\essaytitle}[2]{%
  \title{#1\\[0.3em]\large #2}
  \author{Flyxion\\\normalsize Independent Researcher}
  \date{\today}
}

\pagestyle{plain}

\essaytitle{Sleep Architecture as Record and Commit}{Against a purely hydraulic reading of NREM/REM cycling}
\begin{document}
\maketitle
\begin{abstract}
Adult human sleep is usually described as a sequence of four to six NREM/REM cycles lasting roughly 90--110 minutes, with early-night slow-wave sleep and late-night REM predominance. In one respect that description is plainly correct. In another, it is often framed by an unexamined metaphor: sleep pressure builds and is discharged, as though the architecture of the night were fundamentally a regulated pressure-release system. This essay argues that such language is indispensable for one part of sleep science but insufficient for another. Homeostatic sleep pressure, formalized in Borb\'ely's two-process model as Process S, captures a real monotonic variable. But the content-sensitive work described by active systems consolidation does not look like mere discharge. Slow-wave NREM, especially when slow oscillations, spindles, and hippocampal ripples are coordinated, is better read as a \repair-like interval in which recently encoded traces are reactivated, cross-checked, and selectively rendered \admissible for longer-term cortical integration. REM sleep is then comparatively commit-like: not the whole of consolidation, and not a universally settled mechanism, but a distinct phase in which what has survived prior reactivation is further integrated, especially for procedural and emotional material. Targeted memory reactivation experiments, REM-emotion studies, and the stubborn counterevidence of Process S together show where the hydraulic picture works, where it fails, and why the night's architecture is better understood as alternating continuation passes than as one prolonged discharge event.
\end{abstract}

\section{Introduction}
A healthy adult night of sleep is structured rather than homogeneous. NREM sleep proceeds through lighter and deeper stages, with slow-wave sleep concentrated especially in N3, while REM sleep is marked by rapid eye movements, skeletal muscle atonia, and a low-voltage, mixed-frequency EEG often described as relatively desynchronized compared with deep NREM. Across a typical night, these states alternate in cycles of roughly an hour and a half to a little under two hours, with more slow-wave sleep early and more REM later. What is less often noticed is the picture of what the alternation \emph{is for}. Even outside popular writing, sleep is frequently described by a pressure metaphor: pressure accumulates during wake, deep sleep pays it down, REM either accompanies the process or completes the discharge in a different register.

That metaphor is not invented out of nothing. Borb\'ely's two-process model gave sleep science a durable and fruitful account of homeostatic sleep need and circadian timing. Process S really does rise with waking and decline in sleep, and slow-wave activity is one of its best-known physiological correlates. But once the question shifts from \emph{when and how strongly one sleeps} to \emph{what happens to the content of recent experience during sleep}, a purely hydraulic description begins to lose grip. The active systems consolidation framework associated above all with Jan Born and Susanne Diekelmann describes something more selective: repeated hippocampal-neocortical reactivation during slow-wave sleep gradually stabilizes and redistributes hippocampus-dependent memories into neocortical long-term storage.\footnote{Susanne Diekelmann and Jan Born, ``The memory function of sleep,'' \emph{Nature Reviews Neuroscience} 11, no. 2 (2010). See also Jan Born and Ines Wilhelm, ``System consolidation of memory during sleep,'' \emph{Psychological Research} 76, no. 2 (2012).} Crucially, that physiology is not just loosely synchronous. In human intracranial recordings, slow oscillations, spindles, and hippocampal ripples show hierarchical nesting: spindle activity is biased toward the slow-oscillation up-state, while ripple activity is further concentrated within the spindle phase, giving the coordination claim a real timing structure rather than a hand-waving one.\footnote{Bernhard P. Staresina et al., ``Hierarchical nesting of slow oscillations, spindles and ripples in the human hippocampus during sleep,'' \emph{Nature Neuroscience} 18, no. 11 (2015): 1679--1686.} That is not well described as mere pressure release.

A companion essay in this series has made an analogous anti-hydraulic argument about emotion: the point here is that sleep invites the same correction for different reasons. The proposal is modest. It does not deny sleep homeostasis, and it does not claim that REM cleanly ``finalizes'' every memory while NREM merely ``stores'' it. The claim instead is that the NREM/REM cycle is better read, for memory, as alternating \repair and commit passes on labile traces than as a single release valve operating under two names.

\subsection*{Unexamined premise}
The unexamined premise is that the night's architecture is basically organized by \emph{discharge}: sleep need builds up; NREM and REM are then episodes through which the organism relieves that accumulated load. In the scientific literature this premise is respectable when the topic is homeostatic regulation. In colloquial and semi-technical explanation, however, it often expands beyond its legitimate domain. Because Process S is real, sleep as such gets treated as one large expenditure of a stored quantity, and the alternation of NREM and REM is assimilated to phases of the same release mechanism. On that picture, what matters primarily is how much pressure has been shed, while memory effects are downstream benefits.

The active-systems literature suggests a different priority. If replay during slow-wave sleep is selective, if externally re-presented cues can bias which traces are strengthened, and if REM contributes differently from NREM, then the architecture of the night is not adequately described by a content-blind variable being monotonically drained. The field's often unstated background picture is thus too coarse for the phenomena it now routinely measures.

\subsection*{Temporal and constitutive priority}
Temporal priority and constitutive priority must be separated early. Temporally, the case for homeostatic priority is strong: wakefulness comes first, Process S increases, and sleep follows. Likewise, an episode must be encoded before later sleep can act on it. Nothing in what follows denies that ordinary chronology.

Constitutively, however, the issue is different. What counts as \emph{the same remembered trace continuing} from learning to later recall is not fixed at encoding and then merely preserved. On the active systems view, part of what makes a hippocampus-dependent trace the durable memory it later becomes is precisely that it survives a sleep-dependent pass of reactivation, comparison, and redistribution. The night's work is not merely something that happens \emph{to} an already complete memory object. It helps constitute which traces continue as stable memories at all. That is why the pressure model, though temporally accurate in one dimension, is constitutively incomplete in another.

\section{Minimum apparatus}
Only two pieces of the series vocabulary are needed here.

\begin{definition}
For sleep-dependent memory consolidation, a trace is \admissible when its continued integration with existing cortical structure can proceed without destabilizing the larger memory system that must carry it forward. ``Admissible'' here does not mean true, vivid, or permanent; it means fit to continue through the available consolidation process.
\end{definition}

\begin{claim}
Relative to memory processing, slow-wave NREM is best understood as a \repair-like phase of selective reactivation and cross-checking, while REM is comparatively commit-like: a distinct continuation phase in which what has already survived reactivation is further integrated, recalibrated, or differentially strengthened under a different physiological regime.\footnote{The closer formal match for this specific claim is not the identity material in Part V but the trajectory notation of \emph{Against Chain of Thought: Toward Causally Faithful Oversight via Chain of Memory} (Flyxion), which models reasoning as a sequence of memory states $M_0 \xrightarrow{f_1} M_1 \xrightarrow{f_2} \cdots \xrightarrow{f_k} M_k$ under learned transformations $f_i$, rather than as a linear token trace. Read NREM as one such transformation $f_i$ --- the step that takes a labile, newly encoded $M_i$ and tests it for admissible carry-forward --- and REM as the following $f_{i+1}$, which operates on whatever survived that test rather than on the raw encoding again. The point of borrowing this notation is narrow: it makes explicit that consolidation is a composition of discrete, testable transformations on a state, not a single undifferentiated overnight process. Source: \url{https://github.com/standardgalactic/research-projects/blob/main/compendium/Chain of Memory.txt}.}
\end{claim}

The claim is comparative rather than absolute. NREM does not do only repair; REM does not do only commit. The point is that the alternation is more intelligible as different passes of continuation than as repeated partial drainings of one reservoir.

At this point it helps to distinguish three positions rather than two. A pure hydraulic account, modeled on Process S, says sleep architecture is fundamentally about the rise and fall of a homeostatic pressure variable. Tononi and Cirelli's synaptic homeostasis hypothesis is not identical to that view: it proposes a more specific global downscaling or renormalization of synaptic strength during sleep, roughly proportional to prior waking potentiation, so that plasticity remains sustainable.\footnote{Giulio Tononi and Chiara Cirelli, ``Sleep and the price of plasticity: from synaptic and cellular homeostasis to memory consolidation and integration,'' \emph{Neuron} 81, no. 1 (2014): 12--34.} My claim is a third one. Even if Process S and some form of synaptic renormalization are both real, the record/commit picture is about something neither account settles by itself: which particular traces are replayed, selected, transformed, and carried forward as later-usable memory.

\section{Case one: slow-wave NREM and hippocampal-neocortical dialogue}
Diekelmann and Born's review remains the cleanest statement of the active systems consolidation view. Newly encoded declarative memories are initially dependent on the hippocampus; subsequent sleep, especially slow-wave sleep, supports repeated reactivation of those traces and their gradual redistribution to the neocortex. The physiological picture is not simply that ``deep sleep is good for memory.'' It is that several rhythms are temporally coordinated. Cortical slow oscillations provide alternating down- and up-states; sleep spindles are preferentially nested in the up-states; hippocampal sharp-wave ripples, associated with replay, are coordinated with these events strongly enough to motivate the picture of a structured hippocampal-neocortical dialogue rather than a generic off-line rest state. Staresina et al. sharpened this point in humans by showing hierarchical nesting rather than mere co-occurrence: spindle activity preferentially tracked the slow-oscillation up-state, and ripple events were themselves locked within the spindle cycle.\footnote{Staresina et al., ``Hierarchical nesting of slow oscillations, spindles and ripples in the human hippocampus during sleep.''}

The mechanistic case is also stronger than human EEG and imaging alone. In a classic rodent study, Wilson and McNaughton showed that hippocampal ensemble firing patterns expressed during waking reappeared during subsequent sleep, supplying foundational replay evidence at the level of neuronal populations rather than only macroscopic rhythms.\footnote{Matthew A. Wilson and Bruce L. McNaughton, ``Reactivation of hippocampal ensemble memories during sleep,'' \emph{Science} 265, no. 5172 (1994): 676--679.} More importantly for present purposes, Girardeau et al. went beyond correlation: using real-time ripple-triggered stimulation to selectively suppress hippocampal sharp-wave ripples during sleep, they found impaired later spatial memory, making ripple participation in consolidation a causal claim rather than a merely observational one.\footnote{Gabrielle Girardeau et al., ``Selective suppression of hippocampal ripples impairs spatial memory,'' \emph{Nature Neuroscience} 12, no. 10 (2009): 1222--1223.}

Even here one should resist overclaim. The precise causal weight of each rhythm, and the exact form of their coupling in humans versus animal models, remain active research questions. Some stimulation studies and meta-analytic summaries support the importance of precise slow-oscillation--spindle timing, but debate has not disappeared. Still, the core contrast with the hydraulic model is already visible. A pressure variable, by itself, has no reason to care \emph{which} memory trace is replayed when. The slow-wave consolidation picture does.

That selectivity is why the language of \repair is apt. A trace encoded during waking is not simply stored until its pressure debt is paid. It is re-entered into an existing cortical field of prior traces, schemas, and learned regularities. Some elements cohere well with what is already there; some are weakened; some are perhaps preserved only in transformed form. The result need not be a byte-for-byte restoration of the waking episode. It is enough that a viable continuation be found. The memory system keeps going by integrating recent traces in a way that remains \admissible to the larger organization of cortical storage.

This is also where the active systems account differs from synaptic homeostasis in the narrower Tononi-Cirelli sense already noted above. Global downscaling may be part of what sleep does, but it is still a different proposition from saying that particular hippocampus-dependent traces are selectively replayed and redistributed. A hydraulic picture fits global renormalization more readily than it fits trace-specific replay.

\section{Case two: targeted memory reactivation as evidence for gating}
If sleep merely discharged a scalar quantity, then externally biasing \emph{which} memory is later improved should be difficult to explain. Targeted memory reactivation experiments are therefore especially revealing. In Rasch and colleagues' 2007 study, participants learned object-location associations in the presence of a distinctive odor, and re-exposure to that odor during slow-wave sleep improved later recall; the same cue during wakefulness or REM did not show the same benefit.\footnote{Bj\"orn Rasch, Christian B\"uchel, Steffen Gais, and Jan Born, ``Odor cues during slow-wave sleep prompt declarative memory consolidation,'' \emph{Science} 315 (2007).} Rudoy and colleagues extended the point with auditory cues: sounds associated with particular object-location pairs were replayed during slow-wave sleep in a nap, and later recall improved selectively for the cued items.\footnote{Jonathan D. Rudoy, Joel L. Voss, Christopher E. Westerberg, and Ken A. Paller, ``Strengthening individual memories by reactivating them during sleep,'' \emph{Science} 326 (2009).} The important fact is not only that sleep helps memory, but that experimental intervention can bias the fate of individual traces.

This does not prove that the sleeping brain is literally running an admissions committee. But it does show that the system is not behaving like a blind sink for accumulated need. Cueing during slow-wave sleep appears to tilt the competition among traces: some are preferentially reactivated and later expressed more robustly. That is much closer to a selective admissibility filter than to undifferentiated pressure release.

One might object that the cue simply injects extra activation and the memory gets stronger for that reason alone. Yet even this weaker reading leaves the hydraulic model short of resources. The relevant question becomes why slow-wave sleep offers a privileged window in which externally reinstated context can alter later consolidation outcomes. TMR works not by bypassing the ongoing process, but by entering an already structured period of replay and hippocampal-neocortical exchange.

NREM is not only a favorable chemical milieu. It is a phase in which traces are still open enough to be differentially taken up. ``Record'' is not the perfect word, because the brain is not transcribing a passive log, but compared with REM, slow-wave NREM is where recent traces are most plainly re-presented and checked against broader cortical organization.

\section{Case three: REM sleep as a commit-like phase, with caution}
REM sleep has long resisted simple summary. It is neurochemically distinct from slow-wave NREM, with low noradrenergic tone and relatively cholinergically activated cortical dynamics; it is also the sleep stage most tightly associated with vivid dreaming, muscle atonia, and strong limbic engagement. The memory literature has often linked REM more to procedural learning, associative integration, and emotional memory processing than to the hippocampal systems consolidation story emphasized for slow-wave sleep.

The well-known work from Matthew Walker's group made an influential version of this claim. Studies connected REM physiology, including REM theta activity, with preferential retention or overnight transformation of emotional material, and more broadly proposed a ``sleep to remember, sleep to forget'' picture in which REM helps preserve the informational content of affectively salient experience while reducing its next-day emotional charge.\footnote{For an influential statement of this view, see Els van der Helm and Matthew P. Walker, ``Overnight therapy? The role of sleep in emotional brain processing,'' \emph{Psychological Bulletin} 136, no. 5 (2010), together with related REM-theta work from Walker's lab. The broader literature is mixed on the strength and specificity of REM's role.} Subsequent work has preserved the importance of REM while complicating its exclusivity. Some findings replicate a REM-emotion link; others distribute the work across NREM/REM sequences, find task-dependent effects, or report weaker and less consistent benefits than the stronger narrative suggested.

A harder objection cuts deeper than ``mixed results.'' Strong REM-necessity claims sit uneasily with the well-known observation that pharmacological REM suppression---most famously in patients maintained on monoamine oxidase inhibitor antidepressants, which can suppress REM dramatically and for long periods, often precisely because of an underlying affective affliction requiring treatment---has not yielded the catastrophic emotional-processing or memory deficits that a simple ``no REM, no emotional consolidation'' thesis would predict. That does not prove REM is dispensable. It does show that any strong necessity claim outruns the evidence.

That mixed record matters. If the case for REM were cleanly settled, the present argument would risk simply renaming an accepted sequence. It is better, and more accurate, to say that REM looks \emph{comparatively} commit-like. By the time later-night REM is abundant, the early-night dominance of slow-wave sleep has already done much of the replay-heavy work on newly encoded declarative material. REM then occurs in a state-space where traces have already been partially selected and reorganized. Under its distinct neuromodulatory conditions, integration can proceed differently: associative linking, procedural calibration, and emotional retuning are all plausible candidates for work done on material that is no longer maximally labile in the way it was earlier in the night.

Calling this ``commit'' is therefore a controlled metaphor. It does not mean irreversible finalization or perfect stability. It means that REM is better conceived as a phase that biases continuation \emph{after} a prior pass of replay-dependent selection than as a second kind of pressure release serving the same basic function as NREM.

\section{Disconfirming instance: Process S really does look hydraulic}
The best objection comes from within mainstream sleep regulation itself. Borb\'ely's two-process model and the long line of work around slow-wave activity are not superficial metaphors but robust quantitative achievements. After prolonged wakefulness, slow-wave activity is higher; across a normal night it declines, often roughly exponentially; local waking use can alter local slow-wave expression; sleep deprivation increases the homeostatic drive that subsequent sleep then reduces.\footnote{Alexander A. Borb\'ely, ``A two process model of sleep regulation,'' \emph{Human Neurobiology} 1, no. 3 (1982).} If one asks what variable is most faithfully tracked by the decline of slow-wave activity across the night, ``pressure discharge'' remains a serious answer.

This matters because it identifies a domain in which the hydraulic picture is not only intuitive but explanatory. A content-sensitive account of admissibility and repair does not, by itself, predict the monotonic overnight decline of a homeostatic marker. Nor does it explain why sleep deprivation so reliably boosts subsequent slow-wave intensity even when one is not experimentally interested in declarative consolidation. The present thesis would be empty if it pretended otherwise.

The right conclusion is therefore not replacement but stratification. The pressure model captures one real axis of sleep: the regulation of need, intensity, and timing. The record/commit model captures another: how particular traces survive, are transformed, or fail to continue through the night's architecture. Where the two meet, there may be genuine interaction. Slow-wave sleep is both the strongest expression of Process S and the most prominent arena for replay-dependent consolidation. But the existence of a homeostatic gradient does not collapse the content-sensitive work into mere discharge.

\section{Non-triviality conditions}
The central move would be wrong or empty under at least three conditions.

First, it would fail if the sleep-memory literature could be fully reconstructed from content-blind homeostatic variables alone. TMR results count strongly against that possibility, because they show experimentally manipulable selectivity. Second, the move would fail if sleep science already routinely treated NREM/REM architecture as a repair-and-commit sequence in all but name. It does not. Pieces of the idea are present, but the dominant vocabulary still splits homeostatic regulation, systems consolidation, synaptic renormalization, and emotional processing into adjacent literatures rather than asking what picture of \emph{continuation} unifies them. Third, the move would fail if counterexamples like Process S turned out to absorb all the explanatory burden. They do not; they explain a different kind of regularity.

These conditions are deliberately restrictive. The aim is not to impose external vocabulary where the field has no need of it, but to show that the vocabulary earns its keep precisely where the field's own results become content-selective.

\section{A speculative why}
What follows is tentative. One reason hydraulic framings recur may be that they are exceptionally good at describing variables that are scalar, cumulative, and monotonic. Sleep need often behaves that way. Memory continuation does not.\footnote{A structurally similar point has recently been made about world models used for action: Xun Ma et al., \emph{Rethinking World Models for Safety-Critical Embodied Systems}, arXiv:2609.03774 (2026), \url{https://arxiv.org/abs/2609.03774}, argue that predictive fidelity is the wrong sufficient criterion for a world model used to act, since high likelihood on a modal future can be achieved while discarding low-probability distinctions whose loss is catastrophic for intervention. The present essay makes an analogous claim about sleep: a homeostatic variable can be perfectly well predicted while the content-sensitive question of which traces survive to be usable later is left unanswered.} A remembered episode is not simply a quantity to be spent down; it is a structured object whose persistence depends on whether a later system can carry it forward without destabilizing itself. In that respect sleep may expose a broader truth: a state does not contain enough structure to explain its own continuation, because continuation is partly determined by the regime that later receives, tests, and transforms the state.\footnote{For a formal version of the thought that a state is insufficient to explain its own continuation, see Flyxion, \emph{Contradiction and Continuation}, Part V, where identity-under-transport is treated as an invariant of admissible continuation rather than a primitive re-identification of static states (developed in dialogue with Graham Priest's work on identity and change). Source: \url{https://github.com/standardgalactic/alphabet/blob/core/soup/book-latex-source/main.tex}.}

If that is even partly right, then NREM/REM cycling matters not merely because the organism alternates depths of rest, but because different sleep regimes offer different continuation rules for recent traces. Early-night slow-wave sleep keeps traces open to replay-heavy \repair; later REM biases a different mode of integration. The alternation is functional because no single phase performs the whole job.

\section{Conclusion}
Sleep science has good reason to retain the language of homeostatic pressure. Process S is real, measurable, and explanatory. But once the topic becomes sleep-dependent memory consolidation, a purely hydraulic picture leaves too much out. Slow-wave NREM coordinates replay events that selectively reactivate recent traces and negotiate their cortical future; targeted memory reactivation shows that this negotiation can be biased at the level of particular memories; REM, under a distinct physiological profile, appears to continue that work in a different, comparatively commit-like mode. The point is not to abolish pressure from sleep theory. It is to stop letting a true model of sleep need masquerade as a complete model of sleep architecture. The night is not only a discharge. It is also a sequence of continuation passes by which some traces become fit to last.
\end{document}
